\usepackage[utf8]{inputenc}
\usepackage[T1]{fontenc}
\usepackage[spanish,es-nodecimaldot]{babel}
\usepackage{booktabs,tabularx,amsmath,graphicx}
\usepackage{tikz,listings}
\usetikzlibrary{arrows.meta,positioning,calc,fit,babel}
\usetheme{default}
\setbeamertemplate{navigation symbols}{}
\setbeameroption{hide notes}
\definecolor{azul}{RGB}{20,60,110}
\definecolor{naranja}{RGB}{190,75,20}
\definecolor{verde}{RGB}{25,110,100}
\definecolor{gris}{RGB}{92,103,114}
\setbeamercolor{frametitle}{fg=azul}
\setbeamercolor{structure}{fg=azul}
\setbeamercolor{alerted text}{fg=naranja}
\setbeamercolor{block title}{fg=azul,bg=azul!7}
\setbeamercolor{block body}{bg=azul!3}
\setbeamercolor{block title alerted}{fg=naranja,bg=naranja!10}
\setbeamercolor{block body alerted}{bg=naranja!4}
\setbeamerfont{frametitle}{size=\large,series=\bfseries}
\setbeamerfont{block title}{size=\small,series=\bfseries}
\setbeamersize{text margin left=8mm,text margin right=8mm}
\setbeamertemplate{footline}{%
  \hspace*{8mm}{\scriptsize\color{gris}LC3 · CC531 · Pérez y Pezo}%
  \hfill\hyperlink{fuentes}{\scriptsize\color{azul}Fuentes}\hspace{5mm}%
  \hyperlink{mapa}{\scriptsize\color{azul}Mapa de consultas}%
  \hspace{5mm}{\scriptsize\color{gris}\insertframenumber}\hspace*{8mm}\vspace*{4pt}}
\lstset{
  language=Scala,basicstyle=\ttfamily\scriptsize,
  keywordstyle=\color{azul}\bfseries,commentstyle=\color{gris},
  stringstyle=\color{verde},showstringspaces=false,
  breaklines=true,breakatwhitespace=false,columns=fullflexible,
  keepspaces=true,tabsize=2,aboveskip=1pt,belowskip=1pt,
  literate={á}{{\'a}}1 {é}{{\'e}}1 {í}{{\'{\i}}}1 {ó}{{\'o}}1
           {ú}{{\'u}}1 {ñ}{{\~n}}1 {Á}{{\'A}}1 {É}{{\'E}}1
           {Í}{{\'I}}1 {Ó}{{\'O}}1 {Ú}{{\'U}}1 {Ñ}{{\~N}}1
}
\tikzset{
  paso/.style={draw=azul!35,fill=azul!5,rounded corners=2pt,
    text width=2.8cm,minimum height=0.82cm,align=center,
    inner sep=4pt,font=\scriptsize},
  resultado/.style={paso,draw=verde!60,fill=verde!8},
  apoyo/.style={paso,draw=naranja!50,fill=naranja!6},
  flecha/.style={-{Latex[length=2mm]},draw=azul!65,thick}
}
\newcommand{\fuente}[1]{\par\vspace{3pt}{\scriptsize\color{gris}#1}\par}
\newcommand{\pregunta}[1]{\begin{alertblock}{Pregunta sobre el dataset}#1\end{alertblock}}
\newcommand{\entrada}[1]{\par{\footnotesize\textbf{Entrada y filas.} #1}\par\vspace{3pt}}
\newcommand{\respuesta}[1]{\par\vspace{4pt}{\footnotesize\textbf{Respuesta e interpretación.} #1}\par}
% El código se lee del archivo original para evitar transcribir otra implementación.
\newcommand{\codigo}[3]{\lstinputlisting[firstline=#2,lastline=#3]{../code/#1.scala}}
\newcommand{\origen}[2]{\note{Fuente: \nolinkurl{code/#1/#2.scala}. Salida: \nolinkurl{evidencias/salidas/#2.txt}.}}
\newcommand{\enlace}[2]{\hyperlink{#1}{\textcolor{azul}{\underline{#2}}}}
\newcommand{\flujo}[4]{%
  \par\vspace{2pt}{\centering\begin{tikzpicture}
    \node[paso] (a) at (0,0) {#1};
    \node[paso] (b) at (3.55,0) {#2};
    \node[paso] (c) at (7.10,0) {#3};
    \node[resultado] (d) at (10.65,0) {#4};
    \draw[flecha] (a) -- (b);\draw[flecha] (b) -- (c);\draw[flecha] (c) -- (d);
  \end{tikzpicture}\par}\vspace{3pt}}
% Cada fila distingue la transformación de entrada de su reducción.
\newcommand{\parmr}[4][Map]{%
  \par\vspace{3pt}\noindent\begin{tikzpicture}
    \node[paso,text width=3.0cm,inner sep=3pt] (m)
      {\textbf{1. #1 · #2}\\#3};
    \node[resultado,text width=3.0cm,inner sep=3pt] (r) at (3.85,0)
      {\textbf{2. Reduce · #2}\\#4};
    \draw[flecha] (m) -- (r);
  \end{tikzpicture}\par}
